\begin{tikzpicture}[font=\small,line cap=round,line join=round]\draw[geoTrait](0,0)circle(2.5);\draw[geoTrait](2.5,0)--(-2.5,0) (1.2500,2.1651)--(-1.2500,-2.1651);\fill[otNavy](2.5000,0.0000)circle(1.5pt)node[right,geoLabel]{$ A $};\fill[otNavy](-2.5000,0.0000)circle(1.5pt)node[left,geoLabel]{$ C $};\fill[otNavy](1.2500,2.1651)circle(1.5pt)node[above,geoLabel]{$ B $};\fill[otNavy](-1.2500,-2.1651)circle(1.5pt)node[below,geoLabel]{$ D $};\fill[otNavy](0.0000,0.0000)circle(1.5pt)node[below right,geoLabel]{$ O $};\geoSolution{\draw[geoReponse](2.5,0)--(1.2500,2.1651)--(-2.5,0)--(-1.2500,-2.1651)--cycle;\draw[otGreenDark](2.3500,0.2598)--(2.0902,0.1098)--(2.2402,-0.1500);\draw[otGreenDark](1.4000,1.9053)--(1.1402,1.7553)--(0.9902,2.0151);}\end{tikzpicture}
{\par\smallskip\fontsize{8}{10}\selectfont\color{otMuted}Deux diamètres d’un même cercle ; le rectangle obtenu est tracé en vert au corrigé.\par}
